\section{Regarding NHST}\label{app:nhst}
Null hypothesis significance testing (NHST, \citealt{fisher1925, neymanpearson1933})
% [fisher1925] Fisher (1925), ``Statistical Methods for Research Workers'', Oliver
% and Boyd. Introduced the p-value and the concept of significance testing --- the
% foundational contribution to the p-value side of NHST.
%
% [neymanpearson1933] Neyman & Pearson (1933), ``On the Problem of the Most
% Efficient Tests of Statistical Hypotheses'', Phil. Trans. R. Soc. A 231:289--337.
% Introduced the formal H0/H1 framework, Type I/II errors, and the alpha threshold
% --- the hypothesis-testing side of the hybrid that practitioners use today.
is included here as contextual background rather than as a competing method.
A p-value measures the probability of observing data at least as extreme as those
obtained \emph{assuming the null is true}, $P(\mathrm{data}\mid H_0)$, not the
probability that the null is true given the data, $P(H_0\mid\mathrm{data})$
% [cohen1994] Cohen (1994), ``The earth is round (p < .05)'', American Psychologist
% 49(12):997--1003. The canonical demonstration that p-values cannot be read as
% P(H0|data). Directly supports the claim that NHST carries no null-acceptance
% information.
%
% [wasserstein2016] Wasserstein & Lazar (2016), ``The ASA Statement on p-Values:
% Context, Process, and Purpose'', The American Statistician 70(2):129--133.
% Official ASA position: ``A p-value does not provide a good measure of evidence
% regarding a model or hypothesis.'' Adds broad authoritative institutional backing.
\citep{cohen1994, wasserstein2016};
consequently, NHST can only reject the null or remain inconclusive --- it carries
no mechanism to positively accept it, making it structurally unsuitable for the
precision-and-decisiveness criterion that motivates DPitG.
Understanding its sequential behaviour nevertheless clarifies why the Bayesian
HDI+ROPE framework was developed.

In NHST, at each iteration the experiment stops and $\theta_{\rm null}$ is rejected
whenever the p-value falls below the false positive rate threshold $\alpha_{\rm FPR}$.
Since the decision rule is based solely on this single threshold, data that fail to
trigger rejection yield an inconclusive result rather than evidence in favour of the null.

Using the same fair coin setup ($\theta_{\rm true}=\theta_{\rm null}=0.5$) described
in the Methods section (\ref{sec:methods}), we apply NHST with $\alpha_{\rm FPR}=0.05$
as the stopping and decision criterion. A reliable sequential method should maintain a
low rejection rate throughout when the null hypothesis is true: as demonstrated in
Section~\ref{sec:fair_coin_large_scale}, HDI+ROPE's rejection rate quickly asymptotes near 6\% and
both precision-based methods maintain a zero rejection rate.

Figure~\ref{fig:fair_decisions_nhst} shows the opposite for NHST: the rejection rate
(red solid) grows steadily over time, while the inconclusive proportion (gray dashed)
decreases. To make this trend visible, the simulation is extended to
$N_{\rm max}=30{,}000$ iterations, well beyond the $N_{\rm max}=1{,}500$ used in
the main analysis, since false rejections accumulate slowly and the growing trend
would not be apparent at shorter horizons.

\begin{figure}[h!]
  \centering
  \includegraphics[width=1\textwidth]{fair_experiment_decision_rates_nhst.png}
  \caption{Cumulative decision rates across iterations for $M=50$ fair coin experiments
  ($\theta_{\rm true}=\theta_{\rm null}=0.5$, $\alpha_{\rm FPR}=0.05$,
  $N_{\rm max}=30{,}000$), extending Figure~\ref{fig:fair_decisions} with NHST
  as a fourth panel; it appears in this appendix as NHST lies outside the scope
  of the main analysis.
  At each iteration $N$, the plotted proportions count all experiments that reached
  a decision at or before $N$; their sum is 100\% at every point.
  \textit{Red solid line}: reject $\theta_{\rm null}$.
  \textit{Gray dashed line}: inconclusive.
  The smaller experiment count ($M=50$ vs.\ $2{,}000$) and higher $N_{\rm max}$ reflects the higher
  computational cost of evaluating p-values relative to HDI-based criteria,
  which also accounts for the stepped appearance of the curves.}
  Alt text: A single panel with a logarithmic $x$-axis (iteration,
  $10^{0}$--$3{\times}10^{4}$) and a linear $y$-axis (proportion of $M=50$
  experiments, $0$--$1$).
  Title: $\theta_{\rm true}=0.50$, $\theta_{\rm null}=0.50$,
  $\alpha_{\rm FPR}=0.05$.
  Two curves are shown: a red solid line (reject $\theta_{\rm null}$) and a
  grey dashed line (not reject / inconclusive); at every iteration the two
  proportions sum to $1$.
  A light grey horizontal dot-dashed line marks the $\alpha_{\rm FPR}=0.05$
  reference level near the bottom of the panel.
  At $N=1$ the reject curve opens near $0$ and the inconclusive curve starts
  at $1.0$.
  Around $N\approx10$ the reject curve steps up to $\sim\!0.06$, already near
  the $\alpha_{\rm FPR}$ reference line, then continues rising in irregular
  steps as more experiments accumulate sufficient evidence to reject.
  The inconclusive curve mirrors this descent.
  The two curves cross near $N\approx3{,}000$--$5{,}000$.
  By $N\approx3{\times}10^{4}$ the reject curve reaches $\sim\!0.60$ while
  the inconclusive curve falls to $\sim\!0.40$, illustrating the unbounded
  growth of the false rejection rate under continuous NHST peeking.
  The stepped appearance of both curves reflects the small experiment count
  ($M=50$).
  \label{fig:fair_decisions_nhst}
\end{figure}

As pointed out by \citet{kruschke2015doing}, as evidence accumulates under NHST
the probability of observing a sample extreme enough to trigger rejection grows without
bound --- even when the null hypothesis is true. To paraphrase: ``with infinite
peeking, the false positive rate approaches 100\%.'' This is because the stopping
rule is tied directly to the decision rule with no reference to an effect size: the
more data collected, the more likely an extreme sample will appear. Continuously
monitoring the data and applying a fixed p-value threshold therefore leads to an
inflated false-positive (Type~I) error rate, a manifestation of the multiple-testing
problem
% [simmons2011] Simmons, Nelson & Simonsohn (2011), ``False-Positive Psychology'',
% Psychological Science 22(11):1359--1366. Quantifies that flexible stopping rules
% (``peeking'') inflate false-positive rates even at nominal alpha=0.05 --- the same
% mechanism as continuous sequential NHST. Already cited in the Introduction for this
% claim; repeated here for self-contained reading of the appendix.
\citep{simmons2011}.
